% s6: the fast modes dress; they do not erase (point 4 of the theory). ASCII only.
\section{The fast modes dress; they do not erase}

\spine{4} The fine grain cannot erase the note. A field on the sphere has infinitely many modes finer than the room. My
answer to what they do to the slow note is that they dress it. The logarithmically divergent part matches the running
of the coupling~\cite{GrossWilczek73,Politzer73}; finite effects remain, and controlling them is part of my conjecture.

\textbf{The scale fibration.} $S^3_t$ is itself a group, so fields on it sort by how they transform under left
translations (Peter--Weyl): sectors $j=0,1,2,\dots$, sector $j$ with multiplicity $j+1$. Sector $j=0$ is exactly the
homogeneous fields. So the configuration space is a fibration: a slow \emph{base} (the nine numbers $q$ modulo colour)
and over each point a \emph{fibre} of fast modes, $j\ge1$, locally in a chosen gauge chart. This is not a global
trivialization of the gauge quotient. Here $j$ is twice the representation spin.

\begin{claim}{PROVED}{The slow base is exactly the reduced model}\label{b:pw}
The classical homogeneous restriction and quadratic fluctuation decomposition recover the unsourced reduced
model, with $\mu_{\rm YM}=4\omega^2_{\rm red}$ and background coordinates $M=2gq$ (not the curvature $F$).
The full interacting quantum Hamiltonian also couples slow modes to fast pairs and higher products: it is not a
direct sum of independent $j$-sector theories. A single fast quantum has nonzero left spin and cannot by itself mix
into a left-invariant slow state.
\evidence{Representation theory on $S^3=SU(2)$; the unit map checked to $2\times10^{-8}$~\cite{HelikeReport,ClosedSet}.}
\lines{13}
\end{claim}

On the round sphere van Baal and van den Heuvel keep eighteen degenerate slow modes, left- and
right-invariant~\cite{vanBaalvandenHeuvel94,vanBaal00}; I put the right-invariant ones in the fibre and stretch the
sphere ($t\ne1$). My split rests not on a frequency gap but on the coupling: every slow--fast coupling carries a factor
$g$, through $M=2gq$ near the vacuum chart. That is unconventional, and H2 below must justify it.

\textbf{Born--Oppenheimer.} Freeze the slow variables at $q$. The fast modes form a harmonic system with ground state
$\Omega(q)$ and zero-point energy $E_0(M)=\tfrac12\sum_n\sqrt{\mu_n(M)}$ in the Yang--Mills spectral convention.
Reduced energies are half these energies. On a simply connected chart with a positive fast spectrum and a
nondegenerate normalized real $\Omega$, its Berry connection vanishes locally; this does not assert a global
trivialization across degeneracies. The slow potential receives $\tfrac12E_0(2gq)$ plus further terms of order $g^2$.
The sum diverges, and the regulator must be a function of the spectrum,
never of the sector label: gauge invariance forces a fast shift of exactly $-2$ between the two vacua at $t=1$; a
sector cutoff gives $-8.67$, a heat trace $-1.99999997$.

\begin{claim}{COMPUTED}{The dressing's divergence is asymptotic freedom}\label{b:b0}
$E_0=-\dfrac{b_0}{8\pi^2}\log(\Lambda_{\rm cut}R)\,S_{\rm cl}+E_0^{\rm fin}$ with $b_0=22/3$, the one-loop coefficient
of $SU(2)$. In the saved round-sphere line comparison, the maximum relative discrepancy across seven non-vacuum
points is $7.95\times10^{-4}$; across the five central points it is $6.66\times10^{-5}$.
These are observed discrepancies, not statistical confidence intervals. The proposed subtraction is running $1/g^2$.
\evidence{Heat trace on the Berger slice spectrum~\cite{HelikeReport}; the coefficient of~\cite{GrossWilczek73,Politzer73}.
Computed, not proved; my first pass had the sign wrong, corrected in the register~\cite{ClosedSet}.}
\lines{14}
\end{claim}

This is the place I find most beautiful: the computed ultraviolet logarithm recognizes the known coupling law.
It does not say that the fine grain acts only through that number. Its finite potential and kinetic terms can still
move the note, and asymptotic freedom alone does not prove that the infrared gap survives.

\textbf{Where the fine grain acts.} The frequency equation of \S5 shows how dressing can work without changing the
topology. For an isolated positive-frequency mode, the first-order change is
\begin{equation}\label{eq:dressing}
 \delta\omega=
 \frac{v^T(\delta K_{\rm pot}-\omega^2\delta G_{\rm kin})v}
      {2\omega\,v^TG_{\rm kin}v}.
\end{equation}
For a degenerate block, first diagonalize the corresponding perturbation within that block~\cite{Kato}.
The potential changes the stiffness; the kinetic metric changes the ruler against which that stiffness is measured.
Both can move the note. At a fixed background in the orthonormal reduced coordinates, the potential correction
$\tfrac12E_0^{\rm fin}(2gq)$ gives $\delta K_{\rm pot}=2g^2\nabla_M^2E_0^{\rm fin}$ and hence
$\delta\omega_k=g^2u_k\cdot\nabla_M^2E_0^{\rm fin}\cdot u_k/\omega_k$ for a unit slow mode $u_k$.
This is the potential contribution below; the formula also shows where a kinetic contribution enters.

\begin{claim}{COMPUTED}{The potential part of the dressing widens Helike's lead}\label{b:widen}
At the controls and $g^2=0.002$ the harmonic separation $0.0165$ between Helike and the competitor becomes $0.0191$ (range
$0.0183$--$0.0200$ over regulator windows); at $g^2=0.004$, $0.0216$. This is the potential part of the $O(g^2)$ shift
only.
\evidence{Curvatures of $E_0^{\rm fin}$ along the two modes, $8.89$ and $10.42$, converged to four
digits under the reported spectral-truncation check at fixed prescription, but with regulator-window spreads
$8.65$--$9.08$ and $10.20$--$10.67$ respectively~\cite{HelikeReport,ClosedSet}.}
\lines{15}
\end{claim}

\begin{claim}{OPEN}{The kinetic correction at the same order}\label{b:kinetic}
A second $O(g^2)$ correction is not computed: the non-adiabatic mass of the fast modes dragged by the slow motion (which
screens) plus the Coulomb term of the fast Gauss law (which anti-screens), entering as $-m^{\rm ren}_{u_ku_k}\omega_k$.
Until the full effective Hamiltonian at this order is controlled, the sign of the change in separation is open,
and so is physical H8 (Statement~\ref{b:h8}). Computing this term is necessary, not alone sufficient to settle both.
\evidence{Identified by the referee pass; recorded as the correction to line 15~\cite{ClosedSet}.}
\lines{15, 16}
\end{claim}

Equation~\eqref{eq:dressing} turns this open term into a threshold. In the second-order effective Hamiltonian the
kinetic matrix is $\delta_{ik}-2m^{\rm ren}_{ik}$, so $\delta G_{\rm kin}=2m^{\rm ren}$ to first order and the kinetic
shift of a unit slow mode is $-\omega_k m^{\rm ren}_{u_ku_k}$. At the sourced controls the Helike doublet and the
competitor are separated; a symmetry-preserving correction is treated in their charge blocks, keeping the doublet
together. This is a local expansion with a positive kinetic matrix. For a smooth gapped fast chart with bounded
derivatives, writing the Born--Huang potential as $\Phi(q)=g^2\phi(2gq)$ gives
$\nabla_q^2\Phi=4g^4\nabla_M^2\phi$. Its leading constant cancels in same-well excitation gaps. That local power
counting does not control offsets between different wells or the uniform removal of the fast cutoff.

\begin{claim}{COMPUTED}{How large the kinetic term must be to reverse the order}\label{b:threshold}
Write $m^{\rm ren}=g^2\tilde m$ in the scheme of line~14 ($\zeta$, $\mu R=1$) and
$\kappa=\omega_C\tilde m_{CC}-\omega_H\tilde m_{HH}$ for the competitor and Helike modes. Here $\kappa$ names the
kinetic difference, not the finite-counterterm parameter used elsewhere. At $t=3$, $s=1.15$, combining the recorded
anharmonic coefficient and local quadratic dressing at the same held background gives the linear estimate
\begin{equation}\label{eq:threshold}
 D_{\rm lin}(g^2)=0.01655+g^2\bigl(-2.208+1.258-\kappa\bigr).
\end{equation}
The coefficient $-2.208$ is the reduced model's own anharmonic shift, consistent with the converged separation
$0.01210$ at $g^2=0.002$; $+1.258$ is the fast potential dressing of Statement~\ref{b:widen}.
Within this estimate and for $g^2>0$, Helike lies below this competitor when
$\kappa<0.01655/g^2-0.950$; equality is a tie, and the opposite inequality reverses the order.
The rounded thresholds are $7.33$ at $g^2=0.002$ and $3.19$ at $g^2=0.004$.
This form also handles $0.950+\kappa\le0$, where division by that quantity would mislead.
There is no remainder bound here: these are thresholds for the linear estimate, not certified finite-coupling orderings.

I can also ask how the switch responds. Write $D_0(s)=\omega_C(s)-\omega_H(s)$, and let $a(s)$, $p(s)$ and
$\kappa(s)$ be the anharmonic, potential and kinetic coefficients of this local estimate. If it extends smoothly
to a simple switch, implicit differentiation gives
\begin{equation}\label{eq:dressedswitch}
 \left.\frac{ds^*}{d(g^2)}\right|_{0}
 =-\frac{a(s_0^*)+p(s_0^*)-\kappa(s_0^*)}{D_0'(s_0^*)},\qquad
 D_0'(s_0^*)\simeq2.569 .
\end{equation}
The coefficients belong at the harmonic root $s_0^*=1.14357$, not at $s=1.15$. Freezing them at that nearby control
instead gives $s^*\approx1.1436+(0.950+\kappa)g^2/2.569$, a useful estimate, not the root's evaluated derivative.
Without the fast terms its slope is about $0.86$: at $g^2=0.03,0.05,0.07$ it predicts $1.1693,1.1865,1.2037$,
compared with the registered reduced-model values $1.1674,1.1809,1.1931$ obtained with varying coefficients.
It captures the direction and leading scale, not those values exactly. In the frozen-coefficient estimate the
dressed switch is lower only if $\kappa<1.258$, by $(1.258-\kappa)g^2/2.569$.
\evidence{The saved harmonic frequencies are $1.106657$ and $1.090107$; the potential coefficient is the central
value of the calculation in Statement~\ref{b:widen}, not an uncertainty bound. I have not included displacement
of the equilibrium by the fast force; at fixed sources that can also change the stiffness at $O(g^2)$.
The full matched effective Hamiltonian, a source prescription and remainder control are needed for a physical
switch. The kinetic term remains open.}
\lines{15, 16, 18}
\end{claim}

A script in the repository recomputes every number of this linear estimate from committed receipts
(\texttt{lowest\_note\_threshold.py}).

\begin{claim}{CONJECTURE}{H2 (M3): uniform dressing at small $R$}\label{b:m3}
At $s=1$, on a chart around $q=0$, for $t\ne1$ and small $g(R)$: the levels of the full Hamiltonian (gauge-invariant
spectral cutoff $\Lambda_{\rm cut}$, bare coupling run as above) in a low window converge as $\Lambda_{\rm cut}\to\infty$
and lie within $C\,g^2(R)/R$ of the vacuum-chart reduced levels, uniformly in $\Lambda_{\rm cut}$.
The comparison includes multiplicities and exclusion of extra low states outside the retained sector.
\evidence{Proved pieces: the split, the flat Berry connection, the selection rule. Computed: $b_0$. Templates:
L\"uscher~\cite{Luscher83}, Koller--van Baal~\cite{KollerVanBaal88}, van Baal on $S^3$~\cite{vanBaal00}; rigorous
Born--Oppenheimer for finitely many fast modes with a gap~\cite{Teufel03}.}
\lines{16}
\end{claim}

The route I see is Ba{\l}aban's multiscale method~\cite{Balaban} in Hamiltonian form: integrate the fast modes in
spectral shells; by asymptotic freedom shell $m$ has coupling $g_m=g(2^{-m}R)\le g(R)$, and if each contributes at most
$Kg_m^22^{-m}/R$ after the running of $1/g^2$ (and of a speed of light $c_0$) is subtracted, the sum is geometric and at
most a constant multiple of $g^2(R)/R$. The first step is kinetic consistency, that electric and magnetic couplings run alike, which is
related to Khriplovich's Coulomb-gauge decomposition of $b_0$~\cite{Khriplovich69} and Statement~\ref{b:kinetic}. Granting
H2, well-separated low levels can be compared. An $O(g^2)$ error cannot order channels split only at $O(g^2)$;
that needs H2b's stronger remainder and the physical vacuum chart, not the sourced two-well identification.
But a small room is not the world; the room must now grow.
